\documentclass[12pt,openany,oneside,a4paper,brazil]{abntex2}
%article,beamer,book,,letter,abnTeX2| 

\usepackage[utf8]{inputenc} % Permite digitar acentos normalmente no código
\usepackage[T1]{fontenc}    % Ajusta a codificação das fontes para o PDF
\usepackage[brazil]{babel}

%variavies que serao usada na construção do relatorio
\titulo{Relatório de Estágio}
\autor{Pedro Henrique Peixoto Cardoso}
\data{Outubro de 2026}
\instituicao{Instituto Federal de Educação, Ciência e Tecnologia da Bahia -- Campus Feira de Santana}
\local{Feira de Santana -- BA}
\preambulo{Relatório de Estágio supervisionado apresentado ao curso de Bacharelado em Sistemas de Informação do Instituto Federal de Educação, Ciência e Tecnologia da Bahia -- Campus Feira de Santana, como requisito parcial para avaliação e conclusão da disciplina.}
\orientador{Fernanda}


% --- CUSTOMIZAÇÃO DA CAPA ---
\renewcommand{\imprimircapa}{%
  \begin{capa}%
    \center
    {\ABNTEXchapterfont\large\imprimirinstituicao} % Puxa o nome do IFBA para o topo
    
    \vspace*{\fill}
    
    {\ABNTEXchapterfont\large\imprimirautor}
    
    \vspace*{\fill}
    \begin{center}
    \ABNTEXchapterfont\bfseries\LARGE\imprimirtitulo
    \end{center}
    \vspace*{\fill}
    
    {\large\imprimirlocal}
    \par
    {\large\imprimirdata}
    \vspace*{1cm}
  \end{capa}
}
\makeatletter
\renewcommand{\folhaderostocontent}{
  \begin{center}
    {\ABNTEXchapterfont\large\imprimirautor}
    \vspace*{\fill}\vspace*{\fill}
    \begin{center}
      \ABNTEXchapterfont\bfseries\Large\imprimirtitulo
    \end{center}
    \vspace*{\fill}
    
    \abntex@ifnotempty{\imprimirpreambulo}{%
      \hspace{.45\textwidth}
      \begin{minipage}{.5\textwidth}
        \SingleSpacing
        \imprimirpreambulo
      \end{minipage}%
      \vspace*{\fill}
    }%
    
    % A instrução que imprimia a instituição centralizada foi removida daqui
    
    {\large\imprimirorientadorRotulo~\imprimirorientador\par}
    \abntex@ifnotempty{\imprimircoorientador}{%
      {\large\imprimircoorientadorRotulo~\imprimircoorientador}%
    }%
    \vspace*{\fill}
    
    {\large\imprimirlocal}
    \par
    {\large\imprimirdata}
    \vspace*{1cm}
  \end{center}
}
\makeatother

\begin{document}
%-------------------------------------------------------%
%Elementos agora pre-textuais como capa e contra capa
\pretextual
\imprimircapa
\imprimirfolhaderosto
\newpage
\tableofcontents*
\imprimirorientadorRotulo
\begin{resumo}
\vspace{\onelineskip}
\noindent
\textbf{Palavras-chave}: latex. abntex. editoração de texto.
\end{resumo}

%-------------------------------------------------------%
%ELEMENTOS TEXTUAIS%
\textual

\chapter{Introdução}

\chapter{Desenvolvimento}
\section{teste}

\chapter{Conclusão}
% no section ele tem 2 parametros onde o primeiro é 
%o nome que ficara no sumario usando [] para este caso
%\section*{nome da seção}
%-------------------------------------------------------%
%ELEMENTOS POS-TEXTUAIS
\postextual

\end{document}

% à esquerda	flushleft
% à direita	flushright
% centralizado	center

% estilo	Comando para trecho	Comando para bloco
% monespaçado	\texttt{texto}	{ \ttfamily texto }
% sans-serif	\textsf{texto}	{ \sffamily texto }
% serif	\textrm{texto}	{ \rmfamily texto }
% negrito	\textbf{texto}	{ \bfseries texto }
% itálico	\textit{texto}	{ \itshape texto }
% inclinado	\textsl{texto}	{ \slshape texto }
% maiúsculas	\textsc{texto}	{ \scshape texto }
% sublinhado	\underline{texto}	não tem
% enfatizado	\emph{texto}	não tem

% \tiny
% \scriptsize
% \footnotesize
% \small
% \normalsize
% \large
% \Large
% \LARGE
% \huge
% \Huge

% \vspace{2cm} (adiciona 2 centímetros de espaço vertical).

% \hspace{10pt} (adiciona 10 pontos de espaço horizontal).